\documentclass[english, parskip=full]{scrartcl}
\usepackage[utf8]{inputenc}  % Kodierung der Datei
\usepackage[english]{babel}  % Sprache auf Deutsch 
\usepackage{color}
\usepackage{graphicx, gensymb}
\usepackage{amsfonts, cancel, amssymb}
\usepackage{amsmath, amsthm, array, esint}
\usepackage{booktabs, multirow}  % schönere Tabellen
\usepackage{caption}  % Anpassung der Captions
\usepackage{scrlayer-scrpage}    % Manipulation des Seitenstils
%\pagestyle{scrheadings}  % Stil für die Seite setzen
%\automark{section}  % Abschnittsnamen als Seitenbeschriftung 
%\setheadsepline{.5pt}    % Kopzeile durch Linie abtrennen
\usepackage[hidelinks]{hyperref}  % Links und weitere PDF-Features
\usepackage{typearea, tikz, pgffor, verbatim}
\usetikzlibrary{calc, quotes}
\usepackage[cache=false, outputdir=build]{minted}
\usepackage{placeins} % provides \FloatBarrier

\definecolor{bg}{rgb}{0.9,0.9,0.9}
\newcommand\numberthis{\addtocounter{equation}{1}\tag{\theequation}}
\newcommand{\bra}{}
\newcommand{\kom}{}
\newcommand{\brak}{}
\newcommand{\braket}{}
\newcommand{\intx}{}
\newcommand{\intr}{}
\newcommand{\kla}{}
\newcommand{\klam}{}
\newcommand{\klammer}{}
\newcommand{\oper}{}
\newcommand{\tecxt}{}
\newcommand{\Bra}{}
\newcommand{\Ket}{}
\renewcommand{\i}{\text{i}}
\newcommand{\e}{\text{e}}
\newcommand*{\Fourier}{\textsc{Fourier}\ }
\newcommand*{\F}{\mathcal{F}}
\newcommand*{\sinc}{\text{sinc\,}}
\newcommand{\conv}{\mathop{\scalebox{3.5}{\raisebox{-0.38ex}{\makebox[0.5\width][c]{$\ast$}}}}\limits}

\newcommand*{\titel}{05 Kondition und Stabilität}
\newcommand*{\autor}{Joachim Schwardt und Julian Fleck}

\hypersetup{pdfauthor={\autor}, pdftitle={\titel}}

%\titlehead{titlehead}
\subject{Numerik I - Aufgaben}
\title{\titel}
%\author{\autor}
\author{\begin{tabular}{l|l}
Joachim Schwardt & 4768711 \\ 
Julian Fleck & 4759587\\ 
\end{tabular}}
%\author{\begin{tabular}{|l|l|}
%\hline
%Joachim Schwardt & 4768711 \\ \hline
%Julian Fleck & 4759587\\ \hline
%\end{tabular}}
\date{\today}

\begin{document}

%\begin{titlepage}
\maketitle  % Titel setzen
%\tableofcontents  % Inhaltsverzeichnis setzen
%\end{titlepage}

\section*{Aufgabe 1}
Wir wollen die Fläche $A$ eines Dreiecks mit den Eckpunkten $\vec{r}_j := (x_j,y_j)\in\mathbb{R}^2$ für $j=1,2,3$ berechnen.
Eine Möglichkeit dafür ist $A=\frac{1}{2} |d|$, wobei 
\begin{align}
d &= \det\begin{pmatrix}
x_1 & y_1 & 1 \\ 
x_2 & y_2 & 1 \\
x_3 & y_3 & 1 
\end{pmatrix} = x_1y_2 + x_3y_1 + x_2y_3 - x_3y_2 - x_1y_3 - x_2y_1.
\end{align}
Der Eingabefehler jedes Punktes sei durch $\delta$ in der $\|\cdot\|_\infty$ beschränkt, also
\begin{align}
\|\vec{r}_j - \vec{\tilde{r}}_j \|_\infty &= \max\klammer\left\{ |x_j - \tilde{x}_j|, |y_j - \tilde{y}_j| \right\} \le \delta.
\end{align}
Hierbei seien $\tilde{x}_j, \tilde{y}_j$ jeweils die gestörten Eingaben.
Um eine Abschätzung für den absoluten Fehler zu erhalten, können wir die absolute Kondition verwenden.
Diese ist
\begin{align}
\kappa_\tecxt\text{abs} &= \left\| \begin{pmatrix}
\frac{\partial d}{\partial x_1} & \frac{\partial d}{\partial y_1} & \dots
\end{pmatrix} \right\|_\infty = \max\klammer\left\{ \left\vert \frac{\partial d}{\partial x_1} \right\vert, \left\vert \frac{\partial d}{\partial y_1} \right\vert, \dots \right\} = \\
&= \max\klammer\left\{ |y_2 - y_3|, |x_2-x_3|, \dots \right\}.
\end{align}
Anmerkung: Hier geht irgendetwas schief, weiter unten haben wir einen anderen Ansatz verfolgt. 
Die Abschätzung ist nicht mit em numerischen Experiment kompatibel, aber wir wissen nicht so genau, an welcher Stelle der Fehler steckt... 
\\
Insbesondere gilt dann
\begin{align}
|d - \tilde{d}| &\le \kappa_\tecxt\text{abs} \left\| \begin{pmatrix}
x_1 - \tilde{x}_1 & y_1 - \tilde{y}_1 & \dots
\end{pmatrix}^\top \right\|_\infty = \\
&= \kappa_\tecxt\text{abs} \max \klammer\left\{ |x_1 - \tilde{x}_1|, |y_1 - \tilde{y}_1|, \dots \right\} \le \kappa_\tecxt\text{abs}\delta .
\end{align}
Für die Fläche finden wir daher die Abschätzung
\begin{align}
|A - \tilde{A}| &= \frac{1}{2}\left\vert |d| - |\tilde{d}| \right\vert \le \frac{1}{2}|d - \tilde{d}| \le \frac{1}{2}\kappa_\tecxt\text{abs}\delta.
\end{align}
Der Ausgabefehler skaliert also linear mit dem Eingabefehler.
Zudem ist er proportional zur maximalen Differenz zweier Eckpunkte, was für einen absoluten Fehler zu erwarten war (?).
\\
Auch ohne ``Ausprobieren'' erkennt man aber, dass die Abschätzung nicht für beliebige $\delta$ funktionieren kann.
Sei $y_1=x_2=x_3=y_3=0$:
\begin{align}
d &= x_1y_2,\\
\tilde{d} &= \kla\left( x_1 + \delta_x \right) \kla\left( y_2 + \delta_y \right) = d + y_2\delta_x + x_1\delta_y + \delta_x\delta_y.\\
|d-\tilde{d}| &= |y_2\delta_x + x_1\delta_y + \delta_x\delta_y| %\le |y_2\delta_x + x_1\delta_y| + \delta^2
\end{align}
Für $\delta_x=-\delta$ und $\delta_y=\delta$ ist dann
\begin{align}
|d-\tilde{d}| &= |y_2-x_1 + \delta|\delta,
\end{align}
der Fehler skalliert also wie $\delta^2$, was nicht mit der Abschätzung über die absolute Kondition vereinbar ist.  
Das erklärt auch, warum die Fehlerschranke im numerischen Experiment überschritten wird.
\\ \\
Ein ganz anderer Ansatz nutzt $\vec{1} := (1,1,1)^\top$ und die Abschtzäung von Vektorprodukten über
\begin{align}
A &= \frac{1}{2}|d| = \frac{1}{2} \left\vert \vec{x}\cdot \kla\left( {\vec{y}\times \vec{1}} \right)  \right\vert \le  \frac{1}{2} \|\vec{x}\|_2 \|\vec{y}\|_2 \|\vec{1}\|_2 = \\
&= \frac{1}{2} \sqrt{x_1^2 + x_2^2 + x_3^3} \sqrt{y_1^2 + y_2^2 + y_3^3} \sqrt{3}.
\end{align}
Für $\tilde{x}_j = x_j + \delta_{x_j}$ mit $|\delta_{x_j}| \le \delta$ folgt dann $\|\vec{\delta}_x\|_2 \le \sqrt{3} \delta$, und somit
\begin{align}
|A-\tilde{A}| &\le \frac{1}{2}|d - \tilde{d}|  \frac{1}{2} \left\vert \vec{x}\cdot \kla\left( {\vec{y}\times \vec{1}} \right) -  [\vec{x} + \vec{\delta}_x ]\cdot \kla\left( {[\vec{y} + \vec{\delta}_x]\times \vec{1}} \right) \right\vert = \\
&= \frac{1}{2} \left\vert \vec{\delta}_x\cdot \kla\left( {\vec{\delta}_y\times\vec{1}} \right) + \vec{\delta}_x\kla\left( { \vec{y} \times\vec{1}} \right) + \vec{x}\cdot\kla\left( {\vec{\delta}_y \times\vec{1}} \right) \right\vert \le \\
&=\frac{\sqrt{3}}{2} \kla\left( {\|\vec{x}\|_2 \|\vec{\delta}_y\|_2 + \|\vec{y}\|_2 \|\vec{\delta}_x\|_2 + \|\vec{\delta}_x\|_2 \|\vec{\delta}_y\|_2}  \right) \le \\
&\le \frac{3\delta}{2} \kla\left( {\|\vec{x}\|_2 + \|\vec{y}\|_2 + \sqrt{3}\delta} \right).
\end{align}
Für große $\delta$ deutet sich im numerischen Experiment an, dass diese Abschätzung tatsächlich scharf ist.
Für $\delta\ll |x_j|\dots$ ist in Abbildung \ref{figure:determinant error} allerdings eine große Abweichung zu erkennen.
\begin{figure}[ht!]
\centering
\includegraphics[width=\linewidth]{Determinant_delta_error2.png}
\caption{Die Abschätzung des Fehlers bei der Flächenberechnung von Dreiecken mittels Determinante ist für $N=20000$ Werte von $\delta$ für ein bestimmtes Dreieck dargestellt.
Hierbei sind die Koordinaten $x_j,y_j$ von der Größenordnung 10 ($x_1 = 15$, $y_1 = 25$, $x_2 = 30$, $y_2 = 25$, $x_3 = 20$ und $y_3 = 40$).
Für $\delta \gtrsim 10$ ergibt sich dann wie zu erwarten ein quadratisches Fehlerverhalten.}
\label{figure:determinant error}
\end{figure}
\\
Das quadratische Fehlerverhalten ist aber natürlich nicht besonders interessant, weil die Eingabefehler im Allgemeinen (vermutlich) kleiner als die tatsächlich Eingabe sein werden.
Eigentlich interessiert uns also eher eine Abschätzung unter der Bedingung $\delta \ll |x_j|,\dots$.
Oben hatten wir für die absolute Kondition die $\|\cdot\|_\infty$-Norm angenommen. 
Tatsächlich scheint der richtige Ansatz aber
\begin{align}
\kappa_\tecxt\text{abs} &= \left\| \begin{pmatrix}
\frac{\partial d}{\partial x_1} & \frac{\partial d}{\partial y_1} & \dots
\end{pmatrix} \right\|_1 = \\
&= \left\vert \frac{\partial d}{\partial x_1} \right\vert + \left\vert \frac{\partial d}{\partial y_1} \right\vert + \dots = \\
&= |y_2 - y_3| + |x_2 - x_3| + \dots 
\end{align}
zu sein.
Betrachte zur Rechtfertigung der Verwendung der $\|\cdot\|_1$-Norm einige Extremfälle.
Zunächst sei ausschließlich $\delta_{x_1}\neq 0$, also
\begin{align}
|d - \tilde{d}| &= |\delta_{x_1}||y_2 - y_3| \le \delta |y_2-y_3| \le \delta \kappa_\tecxt\text{abs}.
\end{align}
Sei nun auch $\delta_{x_2}\neq 0$, also
\begin{align}
|d - \tilde{d}| &\le |\delta_{x_1}||y_2 - y_3| + |\delta_{x_2}||y_3 - y_1| \le \delta (|y_2-y_3| + |y_3 - y_1|) \le \delta \kappa_\tecxt\text{abs} .
\end{align}
Hier zeigt sich die Struktur der $\|\cdot\|_1$-Norm für die absolute Kondition.
Die Fehlerabschätzung ist dann wieder
\begin{align}
|A - \tilde{A}| &\le \frac{1}{2}\kappa_\tecxt\text{abs} \delta.
\end{align}
Abbildung \ref{figure:determinant error small delta} bestätigt (bzw. widerlegt zumindest nicht) diese Abschätzung.
\begin{figure}[ht!]
\centering
\includegraphics[width=\linewidth]{Determinant_delta_error3.png}
\caption{Identische Parameter wie in Abbildung \ref{figure:determinant error}, aber für $\delta\ll 10$.
Es ergibt sich das erwartete lineare Fehlerverhalten.}
\label{figure:determinant error small delta}
\end{figure}


%import matplotlib.pyplot as plt
%from matplotlib import special
%special.setup(UseTex=0)
%from numba import njit
%
%@njit
%def get_data(N=10000, dmin=-8, dmax=5):
%    x1 = 15.0
%    y1 = 25.0
%    x2 = 30.0
%    y2 = 25.0
%    x3 = 20.0
%    y3 = 40.0
%    x = np.array([x1, x2, x3])
%    y = np.array([y1, y2, y3])
%    
%    M = np.array([[x1, y1, 1], [x2, y2, 1], [x3, y3, 1]], dtype=np.float64)
%    d = np.linalg.det(M)
%    A = np.abs(d)/2
%    delta_arr = 10.0**np.linspace(dmin, dmax, N)
%    error = np.zeros(N, dtype=np.float64)
%    estimate = np.zeros(N, dtype=np.float64)
%    kappa = np.linalg.norm(np.array([x1 - x2, x1 - x3, x2 - x3, y1 - y2, y1 - y3, y2 - y3]), ord=1)
%
%    xnorm = np.linalg.norm(x, ord=2)
%    ynorm = np.linalg.norm(y, ord=2)
%
%    for i in range(delta_arr.shape[0]):
%        delta = delta_arr[i]
%        x1bar = x1 + np.random.uniform(-delta, delta)
%        x2bar = x2 + np.random.uniform(-delta, delta)
%        x3bar = x3 + np.random.uniform(-delta, delta)
%        y1bar, y2bar, y3bar = np.array([y1, y2, y3]) + np.random.uniform(-delta, delta, 3)
%
%        Mbar = np.array([[x1bar, y1bar, 1], [x2bar, y2bar, 1], [x3bar, y3bar, 1]])
%        dbar = np.linalg.det(Mbar)
%
%        
%        Abar = np.abs(dbar)/2
%        error[i] = np.abs(A - Abar)

%        ### two options: small delta vs. all delta
%        #estimate[i] = kappa * delta / 2
%        estimate[i] = 0.5*np.sqrt(3) *delta * (delta + xnorm + ynorm)
%    return delta_arr, error, estimate
%
%colors = special.Colors()
%
%fig, ax = plt.subplots()
%ax.set_xscale('log')
%ax.set_yscale('log')
%ax.set_xlabel(r"$\delta$")
%ax.set_ylabel(r"$\Delta A$")
%ax.set_xlim(delta[0], delta[-1])
%ax.plot(delta, error, ls='', marker='o', ms=1, mew=1, c=colors.get_color(), label="samples")
%ax.plot(delta, estimate, c=colors.get_color(), label="estimate")
%ax.legend()
%special.polish(fig, ax)


%for delta in 10**np.linspace(-8, 1, 10):
%    x1 = 15
%    y1 = 25
%    x2 = 30
%    y2 = 25
%    x3 = 20
%    y3 = 40
%    M = np.array([[x1, y1, 1], [x2, y2, 1], [x3, y3, 1]])
%    d = np.linalg.det(M)
%    
%    x1bar = x1 + np.random.uniform(-delta, delta)
%    x2bar = x2 + np.random.uniform(-delta, delta)
%    x3bar = x3 + np.random.uniform(-delta, delta)
%    
%    deltay1 = delta - np.abs(x1 - x1bar)
%    y1bar = y1 + np.random.uniform(-deltay1, deltay1)
%    deltay2 = delta - np.abs(x2 - x2bar)
%    y2bar = y2 + np.random.uniform(-deltay2, deltay2)
%    deltay3 = delta - np.abs(x3 - x3bar)
%    y3bar = y3 + np.random.uniform(-deltay3, deltay3)
%    
%    Mbar = np.array([[x1bar, y1bar, 1], [x2bar, y2bar, 1], [x3bar, y3bar, 1]])
%    dbar = np.linalg.det(Mbar)
%    
%    kappa = np.max(np.abs([x1 - x2, x1 - x3, x2 - x3, y1 - y2, y1 - y3, y2 - y3]))
%    A = np.abs(d)/2
%    Abar = np.abs(dbar)/2
%    error = np.abs(A - Abar)
%    estimate = kappa * delta / 2
%        
%    print(f"{delta = :.2e} : {error = :.4e} and {estimate = :.4e}")

 
\newpage
\section*{Aufgabe 2}
Wir betrachten eine Funktion $f : [0,\infty) \rightarrow \mathbb{R}$ mit $f(x) = \sqrt{x+1} - \sqrt{x}$.
\subsection*{a)}
Da $f$ differenzierbar ist, können wir die absolute Kondition über 
\begin{align}
\kappa_\tecxt\text{abs} &= |f'(x)| = \frac{1}{2} \left\vert\frac{1}{\sqrt{x+1}} - \frac{1}{\sqrt{x}} \right\vert
\end{align}
berechnen.
Für große $x$ gilt
\begin{align}
\kappa_\tecxt\text{abs} &\overset{x\,\rightarrow\,\infty }{\longrightarrow} 0,
\end{align}
die Funktion ist dann also gut konditioniert.

\subsection*{b)}
Die Funktion $g : [0,\infty) \rightarrow \mathbb{R}$ mit $g(x) = \kla\left( \sqrt{x+1} + \sqrt{x} \right)^{-1} $.
Dann gilt
\begin{align}
g(x) &= \frac{1}{\sqrt{x+1} + \sqrt{x}} = \frac{\sqrt{x+1} - \sqrt{x}}{\kla\left( \sqrt{x+1} + \sqrt{x} \right)\kla\left( \sqrt{x+1} - \sqrt{x} \right) } = \\
&= \frac{\sqrt{x+1} - \sqrt{x}}{x+1 - x} = \sqrt{x+1} - \sqrt{x} = f(x).
\end{align}


\subsection*{c)}
Für große $x$ ist $f(x)$ zwar gut (absolut) konditioniert, aber bei der Differenzbildung zwischen $\sqrt{x+1}\approx\sqrt{x}$ tritt zunehmend Auslöschung statt.
Bei $g(x)$ werden die Terme allerdings addiert, was numerisch stabil ist.
%Bei der anschließenden Division kann mit voller Genauigkeit gerechnet werden.

\newpage
\section*{Aufgabe 3}
Eine Gleitkommazahl $f$ wird im Computer als
\begin{align}
f &= s\cdot 2^e\cdot \sum_{j=1}^m b_j 2^{-j}
\end{align}
dargestellt, wobei $s\in\{\pm 1\}$ das Vorzeichen angibt, $e$ den Exponenten und die Bits $b_j\in\{0,1\}$ die Mantisse.

\subsection*{a)}
Um eine Dezimalzahl in eine Gleitkommazahl zu überführen, teilen wir diese zunächst in einen ganzzahligen und und fraktionellen Anteil, und dividieren beide wiederholt durch 2.
Der Rest des Ergebnisses ist dabei gerade das Komplement des Bits.
Für die Zahl $125,25$ erhalten wir dabei
\begin{table}[!ht]
\centering
\begin{tabular}{ccccc|c}
Zahl & : & 2 & $=$ & \dots  & Rest\\ \hline
125 & : & 2 & $=$ & 62 & 1\\
62 & : & 2 & $=$ & 31 & 0\\
31 & : & 2 & $=$ & 15 & 1\\
15 & : & 2 & $=$ & 7 & 1\\
7 & : & 2 & $=$ & 3 & 1\\
3 & : & 2 & $=$ & 1 & 1\\
1 & : & 2 & $=$ & 0 & 1\\
\end{tabular}
\end{table}
\\
In Binärdarstellung ist also $125_\tecxt\text{d} = 1111101_\tecxt\text{b}$, wobei die Bits in der Tabelle von unten nach oben ausgelesen werden.
In diesem Fall hätte man auch einfacher darauf kommen können, da $127_\tecxt\text{d} = 1111111_\tecxt\text{b}$, also $125_\tecxt\text{d} = 127_\tecxt\text{d} - 2_\tecxt\text{d} = 1111101_\tecxt\text{b}$.
\\
Für den fraktionellen Anteil multiplizieren wir stattdessen mit 2 modulo 1.
Wir notieren hier als ``Rest'' eine 0 bzw. 1, je nachdem ob das Ergebnis nach der Multiplikation größer als 1 war.
Man findet also
\begin{table}[!ht]
\centering
\begin{tabular}{lcccl|c}
Zahl & $\cdot$ & 2 & $=$ & \dots\ mod 1  & Rest\\ \hline
0,25 & $\cdot$ & 2 & $=$ & 0,5 mod 1 & 0\\
0,5 & $\cdot$ & 2 & $=$ & 0,0 mod 1 & 1\\
\end{tabular}
\end{table}
\\
In Binärdarstellung ist also $0,25_\tecxt\text{d} = 0,01_\tecxt\text{b}$, und damit insgesamt $125,25_\tecxt\text{d} = 1111101,01_\tecxt\text{b}$.
Das können wir auch als
\begin{align}
125,25_\tecxt\text{d} &= 1\cdot 2^6 + 1\cdot 2^5 + 1\cdot 2^4 + 1\cdot 2^3 + 1\cdot 2^2 + 1\cdot 2^0 + 1\cdot 2^{-2} = \\
&= 2^7\cdot\kla\left( 1\cdot 2^{-1} + 1\cdot 2^{-2} + 1\cdot 2^{-3} + 1\cdot 2^{-4} + 1\cdot 2^{-5} + 1\cdot 2^{-7} + 1\cdot 2^{-9}  \right)
\end{align}
schreiben.
Da der führende Term allerdings immer eine 1 sein muss, können wir stattdessen auch
\begin{align}
125,25_\tecxt\text{d} &= 2^6\cdot\kla\left( {\color{blue}{1\cdot 2^{0}}} + 1\cdot 2^{-1} + 1\cdot 2^{-2} + 1\cdot 2^{-3} + 1\cdot 2^{-4} + 1\cdot 2^{-6} + 1\cdot 2^{-8}  \right)
\end{align}
schreiben, aber die blaue $\color{blue}{1}$ nicht speichern.
Es sind also $s=+1$, $e=6$ und $(b_j) = (1,1,1,1,0,1,0,1,0,0,\dots )^\top$.
In IEEE 754 ist eigentlich sogar $e=2^{9} - 1 + 6 = 1000000101_\tecxt\text{b}$.
\\
Dasselbe können wir natürlich auch für die Zahl $0,004150390625$ probieren.
Python kann die Darstellung einfach ausgeben, wir ``schummeln'' aber stattdessen nur ein bisschen, indem wir es für die Divisionen (und den Latex-Code) verwenden.
\begin{table}[!ht]
\centering
\begin{tabular}{lccclc|c}
Zahl & $\cdot$ & 2 & $=$ & \dots & mod 1  & Rest\\ \hline
0.004150390625 & $\cdot$ & 2 & $=$ & 0.00830078125 & mod 1 & 0\\
0.00830078125 & $\cdot$ & 2 & $=$ & 0.0166015625 & mod 1 & 0\\
0.0166015625 & $\cdot$ & 2 & $=$ & 0.033203125 & mod 1 & 0\\
0.033203125 & $\cdot$ & 2 & $=$ & 0.06640625 & mod 1 & 0\\
0.06640625 & $\cdot$ & 2 & $=$ & 0.1328125 & mod 1 & 0\\
0.1328125 & $\cdot$ & 2 & $=$ & 0.265625 & mod 1 & 0\\
0.265625 & $\cdot$ & 2 & $=$ & 0.53125 & mod 1 & 0\\
0.53125 & $\cdot$ & 2 & $=$ & 0.0625 & mod 1 & 1\\
0.0625 & $\cdot$ & 2 & $=$ & 0.125 & mod 1 & 0\\
0.125 & $\cdot$ & 2 & $=$ & 0.25 & mod 1 & 0\\
0.25 & $\cdot$ & 2 & $=$ & 0.5 & mod 1 & 0\\
0.5 & $\cdot$ & 2 & $=$ & 0.0 & mod 1 & 1\\
\end{tabular}
\end{table}
\\
Also ist
\begin{align}
0,004150390625_\tecxt\text{d} &= 2^{-8} + 2^{-12} = 2^{-8}\cdot \kla\left( 2^{0} + 2^{-4} \right),
\end{align}
womit $s=+1$, $e=-8$ und $(b_j) = (0,0,0,1,0,0,\dots)^\top$.



%import struct
%def binary(num):
%    return ''.join('{:0>8b}'.format(c) for c in struct.pack('!f', num))



%def latex_dec_bin(dec, m=52, round_digits=16):
%    f = dec % 1
%    i = int(dec // 1)
%    bits = []
%    while i:
%        r = i % 2
%        bits.append(r)
%        print(f"{i} & : & 2 & " + r"$=$ & " + f"{i // 2} & {r}\\\\")
%        i //= 2
%    bits = bits[::-1]
%    
%    for ctr in range(m - len(bits)):
%        fn = f * 2
%        r = int(fn // 1)
%        bits.append(r)
%        if f:
%            print(f"{f:.{round_digits}f} & $\\cdot$ & 2 & " + r"$=$ & " + f"{fn % 1:.{round_digits}f} & mod 1 & {r}\\\\")
%        f = (2*f) % 1
%
%    return bits

% string = "".join([str(b) for b in bits])

\newpage
\subsection*{b)}
Um zu zeigen, dass $\frac{1}{10}$ nicht exakt in doppelter Genauigkeit darstellbar ist, können wir wieder die ``Multiplikations-Tabelle'' aufstellen. 
%Hier seit man am Python-Output bereits experimentell, dass der Algorithmus tatsächlich alle 52 Bits berechnet, also nie auf die Abbruchbedingung stößt. 
%Insbesondere kann Python die $\frac{1}{10}$ natürlich selbst nicht darstellen, es entstehen also zunehmend Rundungsfehler -- der Latex-Code ist also nicht mehr so einfach zu erzeugen...
\begin{table}[!ht]
\centering
\begin{tabular}{lccclc|c}
Zahl & $\cdot$ & 2 & $=$ & \dots & mod 1  & Rest\\ \hline
0.1 & $\cdot$ & 2 & $=$ & 0.2 & mod 1 & 0\\
0.2 & $\cdot$ & 2 & $=$ & 0.4 & mod 1 & 0\\
0.4 & $\cdot$ & 2 & $=$ & 0.8 & mod 1 & 0\\
\color{blue}0.8 & $\cdot$ & 2 & $=$ & 0.6 & mod 1 & \color{blue}1\\
0.6 & $\cdot$ & 2 & $=$ & 0.2 & mod 1 & 1\\
0.2 & $\cdot$ & 2 & $=$ & 0.4 & mod 1 & 0\\
0.4 & $\cdot$ & 2 & $=$ & 0.8 & mod 1 & 0\\
\color{blue}0.8 & $\cdot$ & 2 & $=$ & 0.6 & mod 1 & \color{blue}1\\
0.6 & $\cdot$ & 2 & $=$ & 0.2 & mod 1 & 1\\
0.2 & $\cdot$ & 2 & $=$ & 0.4 & mod 1 & 0\\
0.4 & $\cdot$ & 2 & $=$ & 0.8 & mod 1 & 0\\
%0.8 & $\cdot$ & 2 & $=$ & 0.6 & mod 1 & 1\\
%0.6 & $\cdot$ & 2 & $=$ & 0.2 & mod 1 & 1\\
%0.2 & $\cdot$ & 2 & $=$ & 0.4 & mod 1 & 0\\
%0.4 & $\cdot$ & 2 & $=$ & 0.8 & mod 1 & 0\\
%0.8 & $\cdot$ & 2 & $=$ & 0.6 & mod 1 & 1\\
%0.6 & $\cdot$ & 2 & $=$ & 0.2 & mod 1 & 1\\
%0.2 & $\cdot$ & 2 & $=$ & 0.4 & mod 1 & 0\\
%0.4 & $\cdot$ & 2 & $=$ & 0.8 & mod 1 & 0\\
%0.8 & $\cdot$ & 2 & $=$ & 0.6 & mod 1 & 1\\
%0.6 & $\cdot$ & 2 & $=$ & 0.2 & mod 1 & 1\\
%0.2 & $\cdot$ & 2 & $=$ & 0.4 & mod 1 & 0\\
%0.4 & $\cdot$ & 2 & $=$ & 0.8 & mod 1 & 0\\
%0.8 & $\cdot$ & 2 & $=$ & 0.6 & mod 1 & 1\\
%0.6 & $\cdot$ & 2 & $=$ & 0.2 & mod 1 & 1\\
%0.2 & $\cdot$ & 2 & $=$ & 0.4 & mod 1 & 0\\
%0.4 & $\cdot$ & 2 & $=$ & 0.8 & mod 1 & 0\\
%0.8 & $\cdot$ & 2 & $=$ & 0.6 & mod 1 & 1\\
\end{tabular}
\end{table}
\\
Nach den ersten drei Nullen wiederholt sich die Tabelle periodisch, also ist $0,1_\tecxt\text{d} = 0,000\overline{1100}_\tecxt\text{b}$.
Genau wie im Dezimalsystem benötigt diese Darstellung aber natürlich unendlich viele Stellen, und damit deutlich mehr als 52.


\subsection*{c)}
Jede Zahl in doppelter Genauigkeit kann mit endlich vielen Stellen als Dezimalzahl dargestellt werden.
Das liegt daran, dass wir nur Linearkombinationen von $b_j\cdot 2^{-j}$ betrachten, wobei $b_j\in\{0,1\}$ natürlich unproblematisch sind.
Die Frage ist also äquivalent dazu, ob die Terme $2^{-j}$ für alle $j=1,\dots,m$ exakt darstellbar sind.
Hierfür gibt es sicherlich viele Beweise, wir versuchen hier einen über Induktion.
Sei $d < 1$ eine Dezimalzahl, sodass 
\begin{align}
d &= \sum_{j=1}^k d_j\cdot 10^{-j},
\end{align}
die Zahl also mit endlich vielen Stellen darstellbar ist, wobei $d_j\in\mathbb{Z}$.
Dann gilt 
\begin{align}
\frac{d}{2} &= \sum_{j=0}^k d_j\cdot 10^{-j-1} \cdot 5 = 5\cdot \sum_{l=1}^{k+1} d_{l-1}\cdot 10^{-l} = 5\cdot \sum_{l=1}^{k+1} \tilde{d}_{l}\cdot 10^{-l} ,
\end{align}
wobei $\tilde{d}_l\in\mathbb{Z}$.
Damit ist $\frac{d}{2}$ also auch exakt darstellbar.
Für $d=\frac{1}{2}=0,5_\tecxt\text{d}$ folgt daraus also induktiv, dass alle $2^{-j}$ für $j=1,\dots,m$ exakt darstellbar sind.


%\begin{thebibliography}{9}
%\bibitem{example} description, \url{https://www.exmapleurl}, day.\,month~2021.
%\end{thebibliography}

\end{document}
